\clearpage
\VerdierPhysicalPageAnchor{0238}%
\begingroup\small\setlength{\parskip}{0pt}\setlength{\emergencystretch}{1.5em}
\setlength{\abovedisplayskip}{4pt plus 2pt minus 2pt}%
\setlength{\belowdisplayskip}{4pt plus 2pt minus 2pt}%

\VerdierSemanticAnchor{ch4d-4-3-4}
\noindent\textbf{4.3.4.}\quad
从 $D(\mathcal A)$ 到 $D(\mathcal A)$ 的谱对象范畴的函子，仍滥用术语称为
\emph{第二典范谱对象函子}。

设 $X$ 是 $\mathcal A$ 中的复形，
$\bigl(X(p,q),\delta(p,q,r)\bigr)$ 是与 $X$ 相联系的第二谱对象。
由 (4.3.2) 和 (1.2.10)，在 $D(\mathcal A)$ 中有典范同构：

\VerdierSemanticAnchor{ch4d-4-3-4-1}
\begin{equation}
X(p,p+1)\xrightarrow{\sim}H^p(X)[-p].
\tag{4.3.4.1}
\end{equation}

\VerdierSemanticAnchor{ch4d-4-3-4-boundedness}
还要注意，当 $X$ 是 $D^*(\mathcal A)$ 的对象
（$*=-,+,b,$ 或“空白”）时，与 $X$ 相联系的第二谱对象事实上取值于
$D^*(\mathcal A)$。

\VerdierSemanticAnchor{ch4d-4-4}
\medskip\noindent\textbf{4.4. 通常的谱序列。}

\VerdierSemanticAnchor{ch4d-4-4-1}
\noindent\textbf{4.4.1.}\quad
设 $\mathcal A$ 是加法范畴，不一定是阿贝尔范畴；设
$F\colon K(\mathcal A)\longrightarrow\mathcal B$ 是到阿贝尔范畴
$\mathcal B$ 的上同调函子，并设
$F^p\colon K(\mathcal A)\longrightarrow\mathcal B$ 是由 $F$ 与
$K(\mathcal A)$ 的平移函子复合所得的函子
（$F^p(X)=F(X[p])$）。此外，设 $Y$ 是 $\mathcal A$ 中的复形。
函子 $F$ 把与 $Y$ 相联系的、取值于 $K(\mathcal A)$ 的第一谱对象
(4.1.4) 变为取值于 $\mathcal B$ 的谱对象。由此得到（第二章 4.3.3）一个
取值于 $\mathcal B$ 的谱序列，它对 $F$ 和 $Y$ 都是函子性的（当后者在
$\operatorname{comp}(\mathcal A)$ 中变化时），称为\emph{函子 $F$ 相对于
复形 $Y$ 的第一谱序列}。对此谱序列，我们始终采用（第二章 4.3.3.1）所引入
的指标方式。

\VerdierSemanticAnchor{ch4d-4-4-2}
\medskip\noindent\textbf{命题 4.4.2.}\quad

\VerdierSemanticAnchor{ch4d-4-4-2-a}
\noindent\textit{a) 函子 $F$ 相对于复形 $Y$ 的第一谱序列的项
$\mathrm I_1^{p,q}$ 由下式给出：}
\[
\mathrm I_1^{p,q}=F^q(Y^p)
\]
\noindent\textit{（这里 $Y^p$ 表示只有一个非零分量的复形；该分量位于
零次，且等于复形 $Y$ 的第 $p$ 个分量）。相应的微分是：}

\VerdierSemanticAnchor{ch4d-4-4-2-a-differential}
\[
\begin{aligned}
d_1&\colon\mathrm I_1^{p,q}\longrightarrow\mathrm I_1^{p+1,q},\\
d_1&=F^q(-d_Y^p)\colon F^q(Y^p)\longrightarrow F^q(Y^{p+1}).
\end{aligned}
\]

\VerdierSemanticAnchor{ch4d-4-4-2-b}
\noindent\textit{b) 因而，同一谱序列的项 $\mathrm I_2^{p,q}$ 可如下描述：
记 $F^{q,\bullet}\colon\operatorname{comp}(K(\mathcal A))
\longrightarrow\operatorname{comp}(\mathcal B)$ 为 $F^q$ 到复形的扩张，
$H_{\mathcal B}^p\colon\operatorname{comp}(\mathcal B)
\longrightarrow\mathcal B$ 为第 $p$ 个上同调对象函子，并记 $Y^\bullet$ 为
$\operatorname{comp}(K(\mathcal A))$ 的下列对象：}

\VerdierSemanticAnchor{ch4d-4-4-2-b-complex}
\[
Y^\bullet=\cdots\longrightarrow Y^{p-1}
\xrightarrow{d_Y^{p-1}}Y^p\xrightarrow{d_Y^p}Y^{p+1}
\longrightarrow\cdots.
\]
\endgroup
